% ATTEMPT_STATUS: unresolved
% RESEARCH_GENERATED: 2026-10-01; GPT-6 Astra Ultra; individual mathematical attempt
% TOP_PROBLEM: {"id": 121, "title": "Sumsets of Perfect Squares", "category_id": 1, "category": {"id": 1, "name": "number_theory", "display_name": "Number Theory", "description": "Properties of integers, prime numbers, Diophantine equations.", "slug": "number-theory", "order_index": 1, "created_at": "2026-07-31T15:26:25.670Z"}, "status": "open"}
% FIRST_POSTED: 2026-10-01
% Imported from: unsolvedmath_additional_500.tex; UnsolvedMath ID 121
% !TeX program = lualatex
% Compile twice: lualatex 121.tex
% Self-contained: no external bibliography, images, or \input files.
% Numeric section labels and visible numbers are original UnsolvedMath IDs.
\documentclass[11pt,a4paper]{article}
\usepackage[margin=24mm,headheight=14pt]{geometry}
\usepackage{fontspec}
\setmainfont{Latin Modern Roman}
\setsansfont{Latin Modern Sans}
\setmonofont{Latin Modern Mono}
\usepackage{amsmath,amssymb,mathtools}
\usepackage{unicode-math}
\setmathfont{Latin Modern Math}
\usepackage{microtype}
\usepackage{enumitem,xurl,needspace}
\usepackage[dvipsnames]{xcolor}
\usepackage{hyperref}
\hypersetup{unicode=true,colorlinks=true,linkcolor=MidnightBlue,urlcolor=MidnightBlue,
 pdftitle={UnsolvedMath Problem 121},pdfauthor={Source-annotated compilation},bookmarksnumbered=true}
\usepackage{bookmark,tocloft,titlesec,fancyhdr}
\setlength{\cftsecnumwidth}{4.2em}
\cftsetpnumwidth{2.5em}
\cftsetrmarg{4em}
\titleformat{\section}{\Large\bfseries\raggedright}{\thesection}{0.7em}{}
\pagestyle{fancy}\fancyhf{}
\fancyhead[L]{\small\sffamily UnsolvedMath research attempt}
\fancyhead[R]{\small Original catalogue IDs}
\fancyfoot[C]{\thepage}
\setlength{\parindent}{0pt}
\setlength{\parskip}{5pt plus 1pt}
\setlength{\emergencystretch}{3em}
\setcounter{tocdepth}{1}\setcounter{secnumdepth}{1}
\setlist[itemize]{leftmargin=3em,itemsep=4pt,topsep=4pt,parsep=0pt}
\allowdisplaybreaks
\newcommand{\N}{\mathbb{N}}
\newcommand{\Z}{\mathbb{Z}}
\newcommand{\Q}{\mathbb{Q}}
\newcommand{\R}{\mathbb{R}}
\newcommand{\C}{\mathbb{C}}
\newcommand{\F}{\mathbb{F}}
\newcommand{\A}{\mathbb{A}}
\newcommand{\PP}{\mathbb{P}}
\newcommand{\E}{\mathbb{E}}
\newcommand{\eps}{\varepsilon}
\newcommand{\dd}{\,\mathrm d}
\newcommand{\sref}[2]{\hyperlink{source:#1:#2}{[S#2]}}
\newif\ifshowcatalogue
\showcataloguetrue
\newcommand{\cataloguescope}[1]{\ifshowcatalogue
 \paragraph{Statement recorded in the supplied source.}
 #1\fi}
\begin{document}
% UnsolvedMath ID 121; source catalogue code GREEN-033

\setcounter{section}{120}

\section{Additive Growth of Sets of Squares}\label{121}

{\small\sffamily UnsolvedMath ID 121 · Number Theory}

\subsection{Definitions and mathematical statement}

\paragraph{Shared notation.}Unless a section says otherwise, $\N=\{1,2,\ldots\}$,
$\Z$ is the ring of integers, $\Q,\R,\C$ are the rational, real, and complex
fields, $[N]=\{1,\ldots,\lfloor N\rfloor\}$, and $\log$ is the natural
logarithm. For real functions of a parameter tending to infinity,
$f\ll g$ and $f=O(g)$ mean $|f|\leq Cg$ eventually for some fixed $C$;
$f\gg g$ reverses this comparison; $f=o(g)$ means $f/g\to0$;
$f\sim g$ means $f/g\to1$; and $f\asymp g$ means both order bounds.
A subscript on $O$ or $\ll$ specifies allowed constant dependence.
The expression $n^{o(1)}$ in an upper bound means growth smaller than
$n^\epsilon$ for each fixed $\epsilon>0$. Local definitions govern any
exception, finite-field convention, or use of the same symbol for another
object. An asymptotic claim is not automatically an effective algorithm.



A square here is a positive integer $u^2$, with $u\in\N$. For a finite set $A$ of such squares, $A+A=\{a+b:a,b\in A\}$ counts distinct sum values, not representations. The conjectured exponent $c>0$ must work uniformly over all finite $A$ with $|A|\geq2$, independently of the magnitudes of its elements. \sref{121}{1} \sref{121}{2}

\paragraph{Statement recorded in the supplied source.}
Is there an absolute constant $c > 0$ such that if $A \subset \mathbb{N}$ is a set of squares of size at least 2, then $|A + A| \geq |A|^{1+c}$? \sref{121}{1}

\paragraph{Residual task reported by the archive.}

The embedded review (2026-08-17) records: Prove a uniform fixed power gain, ideally identify the optimal c, or construct square sets with near-linear doubling that refute it. \sref{121}{1}

\subsection{Short English statement}

Do finite sets of perfect squares always gain a fixed power in size when all pairwise sums are formed?

\subsection{Research attempt}

\paragraph{Provenance and scope.}
This individual attempt was generated by GPT-6 Astra Ultra on 1 October 2026. The method was to derive explicit partial results and test the first proposed extension adversarially. The original statement and sources below are retained. No claim of mathematical novelty or complete solution is made.

\paragraph{Formal target and source scope.}
Let $A=\{b^2:b\in B\}$, where $B$ is a finite set of positive integers, $m=|B|\geq2$, and $M=\max B$. The requested exponent must be independent of both $m$ and the numerical height $M$. The primary problem list makes that uniformity essential. We attack the additive energy through the factorisation of a difference of squares. This yields a useful height-dependent theorem and pinpoints why an apparently near-quadratic estimate is not a solution of the unrestricted problem.

\paragraph{Difference multiplicities and factor pairs.}
For $d>0$, let $r(d)$ count ordered pairs $(a,b)\in B^2$ with $a^2-b^2=d$. Positivity forces $a>b$. The assignment
\[
 (a,b)\longmapsto(u,v)=(a-b,a+b)
\]
is injective and satisfies $uv=d$, $u,v>0$, $u<v$, and $u\equiv v\pmod2$. Conversely such a factor pair recovers $a=(u+v)/2$, $b=(v-u)/2$, although those roots need not lie in $B$. Hence $r(d)\leq\tau(d)$, where $\tau$ is the divisor-counting function. Also $\sum_{d>0}r(d)=m(m-1)/2$.

The additive energy, equivalently counted by equal differences, is
\[
 E(A)=\#\{a_1+a_2=a_3+a_4:a_i\in A\}
     =m^2+2\sum_{d>0}r(d)^2.
\]
The zero difference contributes $m^2$, and positive and negative differences contribute equally. Writing $D(M)=\max_{1\leq d<M^2}\tau(d)$ gives
\[
 E(A)\leq m^2+D(M)m(m-1)\leq(1+D(M))m^2.
\]
Cauchy--Schwarz applied to the sum representation function then proves
\[
 |A+A|\geq\frac{m^4}{E(A)}\geq\frac{m^2}{1+D(M)}.
\]
Every step is valid for arbitrary spacing of the roots; no assumption that $B$ is an interval has entered.

\paragraph{The divisor bound and its dependence.}
For each $\varepsilon>0$ there is $C_\varepsilon$ with $\tau(d)\leq C_\varepsilon d^\varepsilon$. An elementary verification is useful here. For sufficiently large primes $p$, the inequality $a+1\leq p^{\varepsilon a}$ holds for every integer $a\geq0$, because $a+1\leq2^a$ for $a\geq1$. For each of the finitely many smaller primes, the supremum of $(a+1)p^{-\varepsilon a}$ is finite. Multiply those finitely many constants and then multiply over the prime-power factorisation of $d$. This proves the bound with its constant depending only on $\varepsilon$.

Consequently
\[
 |A+A|\geq c_\varepsilon m^2M^{-2\varepsilon}.
\]
If the roots have polynomial height $M\leq m^K$ for a fixed $K$, choose $\varepsilon=1/(4K)$ to obtain $|A+A|\geq c_Km^{3/2}$. In particular a fixed positive power gain is proved in every such height-restricted regime. To remove the multiplicative constant and cover all $m\geq2$, combine this asymptotic inequality with the universal integer-set estimate $|A+A|\geq2m-1$: the latter follows by listing the strictly increasing sums with the smallest element and then with the largest. A sufficiently small exponent depending on $K$ handles the finite range below the asymptotic threshold.

\paragraph{Why this does not prove the catalogue statement.}
The height $M$ is unrestricted in the target. Choosing $\varepsilon$ after seeing $M$ is not legitimate for a uniform fixed exponent, because both $\varepsilon$ and $C_\varepsilon$ then vary with the input. Large common scaling is an additional warning: replacing $B$ by $LB$ preserves $|A+A|$, while it increases $M$ arbitrarily and weakens the displayed estimate. One can divide out a common root factor before applying the method, but primitive root sets can still have height arbitrarily large relative to cardinality. Thus normalisation alone does not remove the obstruction.

A successful energy proof would need an estimate based on how the divisor pairs are distributed within the particular set $B$, rather than the largest divisor count among all integers below $M^2$. The maximum-divisor substitution is the precise place where numerical height replaces combinatorial size.

\paragraph{Verification and exact gap.}
For two distinct squares there are three sums, and the universal $2m-1$ bound is exact. The factor-pair injection excludes $b=0$ consistently with positive roots; allowing zero would require a small separate endpoint adjustment. The energy identity counts ordered pairs and therefore includes both signs with the correct factor two. The result proved here is a complete power-growth theorem under a polynomial-height hypothesis. The unrestricted height-free power gain remains unresolved, and no claim is made that the divisor estimate can supply it without a genuinely stronger input.

\subsection{Sources}

{\small

\begin{itemize}

\item[\hypertarget{source:121:1}{S1}] ULAM.AI, \emph{UnsolvedMath}, supplied \texttt{problems.json}, record 121 (GREEN-033), statement and background. The embedded literature-review date is 2026-08-17; that is the archive’s review date, not a fresh certification. Dataset landing page: \url{https://huggingface.co/datasets/ulamai/UnsolvedMath}.

\item[\hypertarget{source:121:2}{S2}] M.-C. Chang, On problems of Erdős and Rudin, J. Functional Analysis 207 (2004), 444-460, \nolinkurl{doi:10.1016/S0022-1236}(03)00073-9. \url{https://doi.org/10.1016/S0022-1236(03)00073-9}. Bibliographic information imported from the supplied record.

\item[\hypertarget{source:121:3}{S3}] N. Hegyvári, Note on the sumset of squares, arXiv:2504.13230v2 (2025), withdrawal announced. \url{https://arxiv.org/abs/2504.13230}. Bibliographic information imported from the supplied record.

\item[\hypertarget{source:121:4}{S4}] B. Green, 100 Open Problems, Problem 60, version. \url{https://people.maths.ox.ac.uk/greenbj/papers/open-problems.pdf}. The author-hosted document was inspected on 27 September 2026; the archive’s local code is not a problem-number locator in this paper.

\end{itemize}

}

\label{end:121}
\end{document}
